\documentclass[12pt]{article}

\usepackage{filecontents}
%%%%%%%%%%%%%%%%%%%%%%%%%%%%%%%%%%%%%%%%%%%%%%%%%%%%%
%	NOTE:
%					You MUST use sepdflatex (or sepdflatexmk) when you typeset.
%
%%%%%%%%%%%%%%%%%%%%%%%%%%%%%%%%%%%%%%%%%%%%%%%%%%%%%
%  If your document has only "teacher" and "student" versions, then you need not make the changes 
%  described here.
%
%  However, if you want different and/or more versions, you'll need to consider changes to four areas:
%
% (1)  In the bulk of the document (immediately below), use the command \showto to indicate your 
%       audience.
%
%  (2)  Edit the output file names.  See flags below.
%
%  (3)  Edit the names of the audiences.  See flags below.
%
%  (4)  Edit the "current" audience.  See flags below.
%
%%%%%%%%%%%%%%%%%%%%%%%%%%%%%%%%%%%%%%%%%%%%%%%%%%%%%
%%%%%%%%%%%%%%%%%%%%%%%%%%%%%%%%%%%%%%%%%%%%%%%%%%%%%
%%%%%%%%%%%%%%%%%%%%%%%%%%%%%%%%%%%%%%%%%%%%%%%%%%%%%
%%%%%%%%%%%%%%%%%%%%%%%%%%%%%%%%%%%%%%%%%%%%%%%%%%%%%
%%%%%%%%%%%%%%%%%%%%%%%%%%%%%%%%%%%%%%%%%%%%%%%%%%%%%
%%%%%%%%%%%%%%%%%%%%%%%%%%%%%%%%%%%%%%%%%%%%%%%%%%%%%
%%%%%%%%%%%%%%%%%%%%%%%%%%%%%%%%%%%%%%%%%%%%%%%%%%%%%
%%%%%%%%%%%%%%%%%%%%%%%%%%%%%%%%%%%%%%%%%%%%%%%%%%%%%
%
%  Below is the bulk of your document.  Use \showto to indicate audience
%
%%%%%%%%%%%%%%%%%%%%%%%%%%%%%%%%%%%%%%%%%%%%%%%%%%%%%
\begin{filecontents*}{template.tex}
\pagestyle{empty}

\showto{answers}{\hfill ANSWERS\\}
\noindent\rule{7in}{0.4pt}
\begin{center}
{\sc 
\noindent Differential Equations\\
 DEs Practice Problems Set \DEPPSetNum\ 							%CHANGE NUMBER BELOW********
}
\end{center}
\noindent\rule{7in}{0.4pt}

\noindent Submit solutions for extra credit by 1600 \ECDueDate\\
\noindent Corresponding Performance Quiz due by 1600 \PQDueDate\\

\vspace{0.5cm}

%%%%%%%%%%%%%%%%%%%%%%%%%%%%%%%%%%%%%%%%%%%%%%%%%%%%%
\subsection*{Inverse Laplace Transforms}
\begin{DEPP}
\item  (Level 1) Find the inverse Laplace transform of each function.
	\ben[label=(\alph*)]
	\begin{multicols}{2}
	\item $\ds F(s) = \frac{6}{(s-1)^4}$\\
	\item $\ds X(s) = \frac{1}{s^5}$\\
	\item $\ds G(s) = \frac{4}{s^2 + 9}$\\
	\item $\ds Y(s) = \frac{5s}{s^2 + 7}$\\
	\item $\ds R(s) = \frac{3}{(2s+5)^3}$\\
	\item $\ds X(s) = \frac{21}{(s+2)^2 + 16}$\\
	\item $\ds F(s) = \frac{10}{s-10}$\\
	\item $\ds Q(s) = \frac{7s+7}{(s+1)^2 + 12}$\\
	\end{multicols}
	\een
	\vspace{0.25cm}
	
\item  (Level 2) Find the inverse Laplace transform of each function.
	\ben[label=(\alph*)]
	\begin{multicols}{2}
	\item $\ds G(s) = \frac{1}{s^2 + 4s + 8}$\\
	\item $\ds F(s) = \frac{2s+16}{s^2 + 4s + 13}$\\
	\item $\ds Y(s) = \frac{3s-15}{2s^2 - 4s + 8}$\\
	\item $\ds Q(s) = \frac{1}{s^2 + 6s + 9}$\\
	\item $\ds H(s) = \frac{5}{s-6} - \frac{6s}{s^2 + 9} + \frac{3}{2s^2 + 8s + 10}$\\
	\item $\ds X(s) = \frac{7s-4}{s^2 + 36}$\\
	\item $\ds G(s) = \frac{2s-19}{s^2 - 4s+13} + \frac{5}{s-1}$\\
	\item $\ds F(s) = \frac{s}{s^2 + 6s + 11}$\\ 
	\end{multicols}
	\een
	\vspace{0.25cm}
	
\item  (Level 3) Find the inverse Laplace transform of each function.
	\ben[label=(\alph*)]
	\item $\ds X(s) = \frac{6s^2 - 13s + 2}{s^3 - 7s^2 + 6s}$\\
	\item $\ds I(s) = \frac{5s^2 + 34s + 53}{(s+3)^2(s+1)}$\\
	\item $\ds F(s) = \frac{7s^3 - 2s^2 - 3s + 6}{s^3(s-2)}$\\ 
	\een
\newpage

\item  Solve each equation for $F(s),$ then find $\L^{-1}\{ F(s) \}.$
	\ben[label=(\alph*)]
	\item  $\ds s^2 F(s) - 4F(s) = \frac{5}{s+1}$\\
	\item  $ \ds s^2 F(s) + sF(s) - 6F(s) = \frac{s^2 + 4}{s^2 + s}$\\ \vspace{0.5cm}
	\een

\end{DEPP}

\subsection*{Solving DEs with Laplace Transforms}
\begin{DEPP}[resume]
\item  Use Laplace transforms to find the solution $y(t)$ for the initial value problem
	\[y''-2y'+5y = -8e^{-t}, \hspace{0.5cm}y(0) = 2, \hspace{0.5cm} y'(0) = 12.\] 
	\vspace{0.25cm}

\item  Use Laplace transforms to find the solution $y(t)$ for the initial value problem
	\[y'''+4y''+y' -6y = -12, \hspace{0.5cm}y(0) = 1, \hspace{0.5cm} y'(0) = 4, \hspace{0.5cm}y''(0) = -2.\] 
	\vspace{0.25cm}
\end{DEPP}
%%%%%%%%%%%%%%%%%%%%%%%%%%%%%%%%%%%%%%%%%%%%%%%%%%%%%
%%%%%%%%%%%%%%%%%%%%%%%%%%%%%%%%%%%%%%%%%%%%%%%%%%%%%
\showto{answers}{
\newpage
\noindent  Answers:
\begin{DEPP}
\item  \ben[label=(\alph*)]
	\item $\ds f(t) = t^3 e^t$\\
	\item $\ds x(t) = \frac{1}{24}t^4$\\
	\item $\ds g(t) = \frac{4}{3}\sin(3t)$\\
	\item $\ds y(t) = 5 \cos(\sqrt{7}t)$\\
	\item $\ds r(t) = \frac{3}{16}t^2e^{-\frac{5}{2}t}$\\
	\item $\ds x(t) = \frac{21}{4} e^{-2t} \sin(4t)$\\
	\item $\ds f(t) = 10e^{10t}$\\
	\item $\ds q(t) = 7e^{-t}\cos(2\sqrt{3}t)$\\ %\vspace{0.5cm}
	\een
	
\item  \ben[label=(\alph*)]
	\item $\ds g(t) = \frac{1}{2}e^{-2t}\sin(2t)$\\
	\item $\ds f(t) = 2e^{-2t}\cos(3t) + 4e^{-2t}\sin(3t)$\\
	\item $\ds y(t) = \frac{3}{2}e^t\cos(\sqrt{3} t) - 2\sqrt{3}e^t\sin(\sqrt{3} t) $\\
	\item $\ds q(t) = te^{-3t}$\\
	\item $\ds h(t) = 5e^{6t} - 6\cos(3t) + \frac{3}{2}e^{-2t}\sin t$\\
	\item $\ds x(t) = 7\cos(6t) - \frac{2}{3}\sin(6t)$\\
	\item $\ds g(t) = 2e^{2t}\cos(3t) - 5e^{2t}\sin(3t) + 5e^t$\\
	\item $\ds f(t) = e^{-3t}\cos(\sqrt{2} t) - \frac{3}{\sqrt{2}}e^{-3t}\sin(\sqrt{2} t)$\\ %\vspace{0.5cm}
	\een
	
\item  \ben[label=(\alph*)]
	\item $\ds x(t) = e^t + \frac{14}{3}e^{6t} + \frac{1}{3}$\\
	\item $\ds i(t) = -e^{-3t} + 2te^{-3t} + 6e^{-t}$\\
	\item $\ds f(t) = 6e^{2t} + 1 - \frac{3}{2}t^2$\\ %\vspace{0.5cm}
	\een
\newpage
\item  \ben[label=(\alph*)]
	\item  $\ds f(t) = \frac{5}{4}e^{-2t} - \frac{5}{3}e^{-t} + \frac{5}{12}e^{2t}$\\
	\item  $ \ds f(t) = -\frac{13}{30}e^{-3t} + \frac{5}{6}e^{-t} + \frac{4}{15}e^{2t} - \frac{2}{3}$\\ %\vspace{1cm}
	\een

\item  $\ds y(t) = -e^{-t} + 3e^t \cos(2t) + 4e^t \sin(2t)$\\

\item  $\ds y(t) = e^{-3t} - 3e^{-2t} + e^t + 2$

\end{DEPP}
}


%%%%%%%%%%%%%%%%%%%%%%%%%%%%%%%%%%%%%%%%%%%%%%%%%%%%%
\newpage

%\mbox{}\newpage		%If needed

\showto{answers}{\mbox{}\hfill ANSWERS\\}
\noindent\rule{7in}{0.4pt}\\

{\sc 
\noindent Differential Equations 						\showto{student}{\hfill Name:  \line(1,0){200}}\\
DEs Performance Quiz  \DEPPSetNum\					\showto{student}{\hfill HR:  \line(1,0){200}}\\
Due by 1600 \PQDueDate  \showto{student}{\hfill Circle your section: \hspace{0.73in}\mbox{}\\
														\mbox{} \hfill MWF1000
															\hspace{0.5cm}MWF1300 \hspace{0.5cm}TR0925
															\hspace{0.5cm}TR1050}\\
}
\noindent\rule{7in}{0.4pt}
\vspace{0.5cm}

\noindent  For this Performance Quiz, you are permitted to use your own notes/work and your
textbook, but you are NOT AUTHORIZED to seek help from anyone
other than MAJ Libertini or MAJ Bliss.\\

%%%%%%%%%%%%%%%%%%%%%%%%%%%%%%%%%%%%%%%%%%%%%%%%%%%%%
\begin{PQ}
\item    Find the inverse Laplace transform of  $W(s) = \ds \frac{-2s+6}{s^2 + 4}.$\\ \vspace{2in}

\item    Evaluate \hspace{0.25cm}$\ds \L^{-1}\left\{ \frac{s^2 + 6s + 9}{(s+4)(s-1)(s-2)} \right\}$.\\ \label{Q2}

	{\it  Note:  If you've answered correctly, some of the coefficients will be fractions.}\\ \vfill
	
\begin{center} (There's more on the back!)
\end{center}
\newpage	
\item    USE LAPLACE TRANSFORMS to solve the following IVP.
	\[ y'' - 3y' + 2y = e^{-4t}, \hspace{0.5cm} y(0) = 1, \hspace{0.25cm} y'(0) = 5. \]
	
	{\it  Hint:  Your work from question \ref{Q2} should be helpful in answering this question!}\\ \vfill
	
\end{PQ}
%%%%%%%%%%%%%%%%%%%%%%%%%%%%%%%%%%%%%%%%%%%%%%%%%%%%%
%%%%%%%%%%%%%%%%%%%%%%%%%%%%%%%%%%%%%%%%%%%%%%%%%%%%%
\showto{answers}{
\newpage
\noindent  Answers:
\begin{PQ}
\item    $\ds w(t) = -2\cos(2t) + 3\sin(2t)$\\

\item    $\ds \frac{1}{30}e^{-4t} - \frac{16}{5}e^t + \frac{25}{6}e^{2t}$\\
	
\item    $\ds y(t) = \frac{1}{30}e^{-4t} - \frac{16}{5}e^t + \frac{25}{6}e^{2t}$\\
	
\end{PQ}
}
%%%%%%%%%%%%%%%%%%%%%%%%%%%%%%%%%%%%%%%%%%%%%%%%%%%%%
%%%%%%%%%%%%%%%%%%%%%%%%%%%%%%%%%%%%%%%%%%%%%%%%%%%%%
\end{filecontents*}
%%%%%%%%%%%%%%%%%%%%%%%%%%%%%%%%%%%%%%%%%%%%%%%%%%%%%
%%%%%%%%%%%%%%%%%%%%%%%%%%%%%%%%%%%%%%%%%%%%%%%%%%%%%
%%%%%%%%%%%%%%%%%%%%%%%%%%%%%%%%%%%%%%%%%%%%%%%%%%%%%
%%%%%%%%%%%%%%%%%%%%%%%%%%%%%%%%%%%%%%%%%%%%%%%%%%%%%
%%%%%%%%%%%%%%%%%%%%%%%%%%%%%%%%%%%%%%%%%%%%%%%%%%%%%
%%%%%%%%%%%%%%%%%%%%%%%%%%%%%%%%%%%%%%%%%%%%%%%%%%%%%
%%%%%%%%%%%%%%%%%%%%%%%%%%%%%%%%%%%%%%%%%%%%%%%%%%%%%
%%%%%%%%%%%%%%%%%%%%%%%%%%%%%%%%%%%%%%%%%%%%%%%%%%%%%
%%%%%%%%%%%%%%%%%%%%%%%%%%%%%%%%%%%%%%%%%%%%%%%%%%%%%
%  This stuff tells LaTeX to produce multiple versions.
%%%%%%%%%%%%%%%%%%%%%%%%%%%%%%%%%%%%%%%%%%%%%%%%%%%%%
\ifx\conditionmacro\undefined
  \immediate\write18{%
    pdflatex --jobname="answers_\jobname"			%EDIT (2) HERE AS NEEDED
    "\gdef\string\conditionmacro{1}\string\input\space\jobname"  
  }%
  \immediate\write18{%
    pdflatex --jobname="\jobname"						%EDIT (2) HERE AS NEEDED
    "\gdef\string\conditionmacro{2}\string\input\space\jobname"
 }%
%  \immediate\write18{%
%    pdflatex --jobname="solutions_\jobname"
%    "\gdef\string\conditionmacro{3}\string\input\space\jobname"
%  }%
  \expandafter\stop
\fi
%%%%%%%%%%%%%%%%%%%%%%%%%%%%%%%%%%%%%%%%%%%%%%%%%%%%%
%  Here's the preamble.
%%%%%%%%%%%%%%%%%%%%%%%%%%%%%%%%%%%%%%%%%%%%%%%%%%%%%
%\usepackage[cm]{fullpage}
\setlength{\topmargin}{-1in}
\setlength{\oddsidemargin}{-0.25in}
\setlength{\evensidemargin}{-0.25in}
\setlength{\textwidth}{7in}
\setlength{\textheight}{10in}

\usepackage{enumitem}
\usepackage{hyperref}
\usepackage{amssymb}
\usepackage{amsmath}
\usepackage{multicol}
\usepackage{array, latexsym}
\usepackage{times}
\usepackage{time}
\usepackage{subfigure}
\usepackage{graphics}
\usepackage{graphicx}
\usepackage{epstopdf}
\usepackage[T1]{fontenc}
\newcommand{\To}{\Rightarrow}
\newcommand{\del}{\nabla}
\newcommand{\R}{\mathbb{R}}
\newcommand{\C}{\mathcal{C}}
\newcommand{\N}{\mathcal{N}}
\newcommand{\A}{\mathcal{A}}
\renewcommand{\O}{\mathcal{O}}
% \renewcommand{\N}{\mathbb{N}}
\newcommand{\ben}{\begin{enumerate}}
\newcommand{\een}{\end{enumerate}}
\newcommand{\eps}{\varepsilon}
\renewcommand{\epsilon}{\varepsilon}
\newcommand{\p}{\partial}
\newcommand{\<}{\langle}
\renewcommand{\>}{\rangle}
\renewcommand{\subset}{\subseteq}
\newcommand{\cont}{\Rightarrow \Leftarrow}
\newcommand{\back}{\backslash}
\newcommand{\norm}[1]{\|{#1}\|}
\newcommand{\abs}[1]{|{#1}|}
\newcommand{\ip}[1]{\langle{#1}\rangle}
\newcommand{\m}[1]{{\mathbf{#1}}}
\newcommand{\sts}{\m{S}^\textrm{T}\m{S}}
\newcommand{\rank}{\textrm{rank}}
\newcommand{\X}{\mathcal{X}}
\newcommand{\Y}{\mathcal{Y}}
\usepackage{mathrsfs}
\renewcommand{\L}{\mathscr{L}}
\def\ds{\displaystyle}
\def\d{\partial}
\renewcommand{\arraystretch}{1.5}


\newcommand\DEPPSetNum{10}									%CHANGE NUMBER HERE, TOO********
\newcommand\ECDueDate{Friday, 31 March}
\newcommand\PQDueDate{Wednesday, 5 April}

\newlist{DEPP}{enumerate}{1}
\setlist[DEPP]{label = \DEPPSetNum.\arabic*.}

\newlist{PQ}{enumerate}{1}
\setlist[PQ]{label = PQ\DEPPSetNum.\arabic*.}											

%%%%%%%%%%%%%%%%%%%%%%%%%%%%%%%%%%%%%%%%%%%%%%%%%%%%%
%  Here we tell LaTeX the names of the various audience categories it should be prepared for.
%  Add more and/or edit as needed.
%%%%%%%%%%%%%%%%%%%%%%%%%%%%%%%%%%%%%%%%%%%%%%%%%%%%%
\usepackage{multiaudience}
\SetNewAudience{answers}									%EDIT (3) HERE AS NEEDED
\SetNewAudience{student}									%EDIT (3) HERE AS NEEDED
%\SetNewAudience{solutions}


%%%%%%%%%%%%%%%%%%%%%%%%%%%%%%%%%%%%%%%%%%%%%%%%%%%%%
\begin{document}

%%%%%%%%%%%%%%%%%%%%%%%%%%%%%%%%%%%%%%%%%%%%%%%%%%%%%
%  Here's where we tell LaTeX what needs to go into each document.  First, we tell it which 
%  audience it should use, and then we send it up to the "bulk" of the document above.
%%%%%%%%%%%%%%%%%%%%%%%%%%%%%%%%%%%%%%%%%%%%%%%%%%%%%
\ifnum\conditionmacro=1 
	\def\CurrentAudience{answers}							%EDIT (4) HERE AS NEEDED
	\pagestyle{empty}

\showto{answers}{\hfill ANSWERS\\}
\noindent\rule{7in}{0.4pt}
\begin{center}
{\sc
\noindent Differential Equations\\
 DEs Practice Problems Set \DEPPSetNum\        %CHANGE NUMBER BELOW********
}
\end{center}
\noindent\rule{7in}{0.4pt}

\noindent Submit solutions for extra credit by 1600 \ECDueDate\\
\noindent Corresponding Performance Quiz due by 1600 \PQDueDate\\

\vspace{0.5cm}

%intro to LTs / basic LTs.
%%%%%%%%%%%%%%%%%%%%%%%%%%%%%%%%%%%%%%%%%%%%%%%%%%%%%
\subsection*{The Definition of Laplace Transform}

\noindent  As a reminder, the Laplace Transform of function $f(t)$ is a function $F(s)$ as follows:
 \[ F(s) = \L\{ f(t) \} = \int_0^{\infty}e^{-st}f(t) dt, \]
 provided this integral is finite.\\

\begin{DEPP}
\item    Use the definition of Laplace Transform (above) to find the Laplace Transform of $h(t) = 7t.$
 Include appropriate restrictions on $s$ to guarantee the convergence of the integral.\\

\item    Use the definition of Laplace Transform (above) to find the Laplace Transform of $w(t) = e^{kt},$
 where $k$ is some constant.  Include appropriate restrictions on $s$ to guarantee the
 convergence of the integral.\\

\item   Use the definition of Laplace Transform (above) to find the Laplace Transform of $y(t) = 7t + e^{5t}.$
  Include appropriate restrictions on $s$ to guarantee the convergence of the integral.\\

\item  Use the definition of Laplace Transform (above) to find the Laplace Transform of  the function
 \[d(t) = \left\{ \begin{array}{ll} 0, & t<0\\
       7, & 0 \le t < 3\\
       0, & t \ge 3\\
       \end{array} \right.\]  %= 7U(t) - 7U(t-3)
 Include appropriate restrictions on $s$ to guarantee the convergence of the integral.\\

 {\it Note:  This is a really important skill!  If we can take the Laplace Transform of
 piecewise-defined functions like this one, then we can solve differential equations that
 we couldn't solve before!)}\\
%\item    $\ds H(s) = \frac{7}{s^2},$ $s>0$\\
%
%\item    $\ds W(s) = \frac{1}{s-k},$  $s > k$ \\
%
%\item   $\ds Y(s) = \frac{2}{s^2} + \frac{1}{s-5},$  $s>5$

%\item  $\ds D(s) = \frac{7}{s} - \frac{7e^{-3s}}{s}$
\end{DEPP}
\newpage
\subsection*{Using Table to Find LTs}
\begin{DEPP}[resume]
\item  Use the table and properties provided to determine the Laplace transform of each function.
Include resulting restrictions on $s.$
 \ben[label=(\alph*)]
 \item  $\ds f(t) = e^{-2t}\sin(2t) + t^2 e^{3t}$\\
 \item  $\ds g(t) = t^3 e^{-9t}$\\
 \item  $\ds h(t) = e^{2t} - t^3 + t^2 - \sin (5t)$\\
 \item  $\ds q(t) = 8t\cos(6t) + e^{3t}\sin(4t)$\\
 \item  $\ds f(t) = t^2 \sin(3t)$\\
 \een
%\item  \ben [label=(\alph*)]
% \item  $\ds F(s) = \frac{2}{(s+2)^2 + 4} + \frac{2}{(s-3)^3}, \hspace{0.25cm}s>3$\\
% \item  $\ds G(s) = \frac{6}{(s+9)^4}, \hspace{0.25cm}s>-9$\\
% \item  $\ds H(s) = \frac{1}{s-2} - \frac{6}{s^4} + \frac{2}{s^3} - \frac{5}{s^2 + 25},\hspace{0.25cm}s>2$\\
% \item  $\ds Q(s) = \frac{8(s^2-36)}{(s^2 + 36)^2} + \frac{4}{(s-3)^2 + 16}, \hspace{0.25cm}s>3$\\
% \item  $\ds F(s) = \frac{18(s^2 - 3)}{(s^2+9)^3},\hspace{0.25cm}s>0$\\
% \een
\end{DEPP}
%%%%%%%%%%%%%%%%%%%%%%%%%%%%%%%%%%%%%%%%%%%%%%%%%%%%%
%%%%%%%%%%%%%%%%%%%%%%%%%%%%%%%%%%%%%%%%%%%%%%%%%%%%%
\showto{answers}{
\newpage
\noindent  Answers:
\begin{DEPP}
\item    $\ds H(s) = \frac{7}{s^2},$ $s>0$\\

\item    $\ds W(s) = \frac{1}{s-k},$  $s > k$ \\

\item   $\ds Y(s) = \frac{2}{s^2} + \frac{1}{s-5},$  $s>5$\\

\item  $\ds D(s) = \frac{7}{s} - \frac{7e^{-3s}}{s},$  $s>0$\\

\item  \ben [label=(\alph*)]
 \item  $\ds F(s) = \frac{2}{(s+2)^2 + 4} + \frac{2}{(s-3)^3}, \hspace{0.25cm}s>3$\\
 \item  $\ds G(s) = \frac{6}{(s+9)^4}, \hspace{0.25cm}s>-9$\\
 \item  $\ds H(s) = \frac{1}{s-2} - \frac{6}{s^4} + \frac{2}{s^3} - \frac{5}{s^2 + 25},\hspace{0.25cm}s>2$\\
 \item  $\ds Q(s) = \frac{8(s^2-36)}{(s^2 + 36)^2} + \frac{4}{(s-3)^2 + 16}, \hspace{0.25cm}s>3$\\
 \item  $\ds F(s) = \frac{18(s^2 - 3)}{(s^2+9)^3},\hspace{0.25cm}s>0$\\
 \een
\end{DEPP}
}


%%%%%%%%%%%%%%%%%%%%%%%%%%%%%%%%%%%%%%%%%%%%%%%%%%%%%
\newpage

%\mbox{}\newpage  %If needed

\showto{answers}{\mbox{}\hfill ANSWERS\\}
\noindent\rule{7in}{0.4pt}\\

{\sc
\noindent Differential Equations       \showto{student}{\hfill Name:  \line(1,0){200}}\\
DEs Performance Quiz  \DEPPSetNum\   \showto{student}{\hfill HR:  \line(1,0){200}}\\
Due by 1600 \PQDueDate \showto{student}{\hfill Circle your section: \hspace{0.73in}\mbox{}\\
              \mbox{} \hfill MWF1000
               \hspace{0.5cm}MWF1300 \hspace{0.5cm}TR0925
               \hspace{0.5cm}TR1050}\\
}

\noindent\rule{7in}{0.4pt}
\vspace{0.5cm}

\noindent  For this Performance Quiz, you are permitted to use your own notes/work and your
textbook, but you are NOT AUTHORIZED to seek help from anyone
other than MAJ Libertini or MAJ Bliss.\\

%%%%%%%%%%%%%%%%%%%%%%%%%%%%%%%%%%%%%%%%%%%%%%%%%%%%%
\begin{PQ}
\item   {\bf Use the definition of Laplace Transform} to find the Laplace Transform of

  \[y(t) = \left\{ \begin{array}{ll} 0, & t<4\\
        6, & t \ge 4\\
       \end{array} \right.\]    %=6U(t-4)

  Include appropriate restrictions on $s$ to guarantee the convergence of the integral.\\

  {\it Note:  Credit will only be assigned for work that demonstrates how the definition was
  used to find the LT.}\\
\vfill
\begin{center}{\it (There's more on the back!)}
\end{center}
\newpage
\item  Use the table and properties provided to determine the Laplace transform of each function.
Include resulting restrictions on $s.$
 \ben[label=(\alph*)]
 \item  $\ds g(t) = 17t^3 - 9te^{-4t}$\\ \vfill
 \item  $\ds h(t) = t^2 \cos(8t)$\\ \vfill
 \een
\end{PQ}
%%%%%%%%%%%%%%%%%%%%%%%%%%%%%%%%%%%%%%%%%%%%%%%%%%%%%
%%%%%%%%%%%%%%%%%%%%%%%%%%%%%%%%%%%%%%%%%%%%%%%%%%%%%
\showto{answers}{
\newpage
\noindent  Answers:
\begin{PQ}
\item  $\ds Y(s) = \frac{6e^{-4s}}{s},$  $s>0$

\item  \ben[label=(\alph*)]
 \item  $\ds G(s) = \frac{102}{s^4} - \frac{9}{(s+4)^2},$ $s>0$\\

 \item  $\ds H(s) = -16 \cdot \frac{64-3s^2}{(s^2+64)^3}$
  %H(s) = (-1)^2 *d^2/ds^2[8/(s^2 + 64)]
  %  = d/ds[( (s^2+64)(0) - 8(2s)) / (s^2+64)^2]
  %  = -16*d/ds[s / (s^2+64)^2]
  %  = -16 * [ ( (s^2+64)^2 - s*2(s^2+64)(2s) ) / (s^2+64)^4]
  %  = -16 * [ ( (s^2 + 64) - 4s^2 ) / (s^2+64)^3]
  %  = -16 * [ ( -3s^2 + 64) / (s^2+64)^3]
 \een
\end{PQ}
}
%%%%%%%%%%%%%%%%%%%%%%%%%%%%%%%%%%%%%%%%%%%%%%%%%%%%%
%%%%%%%%%%%%%%%%%%%%%%%%%%%%%%%%%%%%%%%%%%%%%%%%%%%%%

\fi 

\ifnum\conditionmacro=2 
	\def\CurrentAudience{student}							%EDIT (4) HERE AS NEEDED
	\pagestyle{empty}

\showto{answers}{\hfill ANSWERS\\}
\noindent\rule{7in}{0.4pt}
\begin{center}
{\sc
\noindent Differential Equations\\
 DEs Practice Problems Set \DEPPSetNum\        %CHANGE NUMBER BELOW********
}
\end{center}
\noindent\rule{7in}{0.4pt}

\noindent Submit solutions for extra credit by 1600 \ECDueDate\\
\noindent Corresponding Performance Quiz due by 1600 \PQDueDate\\

\vspace{0.5cm}

%intro to LTs / basic LTs.
%%%%%%%%%%%%%%%%%%%%%%%%%%%%%%%%%%%%%%%%%%%%%%%%%%%%%
\subsection*{The Definition of Laplace Transform}

\noindent  As a reminder, the Laplace Transform of function $f(t)$ is a function $F(s)$ as follows:
 \[ F(s) = \L\{ f(t) \} = \int_0^{\infty}e^{-st}f(t) dt, \]
 provided this integral is finite.\\

\begin{DEPP}
\item    Use the definition of Laplace Transform (above) to find the Laplace Transform of $h(t) = 7t.$
 Include appropriate restrictions on $s$ to guarantee the convergence of the integral.\\

\item    Use the definition of Laplace Transform (above) to find the Laplace Transform of $w(t) = e^{kt},$
 where $k$ is some constant.  Include appropriate restrictions on $s$ to guarantee the
 convergence of the integral.\\

\item   Use the definition of Laplace Transform (above) to find the Laplace Transform of $y(t) = 7t + e^{5t}.$
  Include appropriate restrictions on $s$ to guarantee the convergence of the integral.\\

\item  Use the definition of Laplace Transform (above) to find the Laplace Transform of  the function
 \[d(t) = \left\{ \begin{array}{ll} 0, & t<0\\
       7, & 0 \le t < 3\\
       0, & t \ge 3\\
       \end{array} \right.\]  %= 7U(t) - 7U(t-3)
 Include appropriate restrictions on $s$ to guarantee the convergence of the integral.\\

 {\it Note:  This is a really important skill!  If we can take the Laplace Transform of
 piecewise-defined functions like this one, then we can solve differential equations that
 we couldn't solve before!)}\\
%\item    $\ds H(s) = \frac{7}{s^2},$ $s>0$\\
%
%\item    $\ds W(s) = \frac{1}{s-k},$  $s > k$ \\
%
%\item   $\ds Y(s) = \frac{2}{s^2} + \frac{1}{s-5},$  $s>5$

%\item  $\ds D(s) = \frac{7}{s} - \frac{7e^{-3s}}{s}$
\end{DEPP}
\newpage
\subsection*{Using Table to Find LTs}
\begin{DEPP}[resume]
\item  Use the table and properties provided to determine the Laplace transform of each function.
Include resulting restrictions on $s.$
 \ben[label=(\alph*)]
 \item  $\ds f(t) = e^{-2t}\sin(2t) + t^2 e^{3t}$\\
 \item  $\ds g(t) = t^3 e^{-9t}$\\
 \item  $\ds h(t) = e^{2t} - t^3 + t^2 - \sin (5t)$\\
 \item  $\ds q(t) = 8t\cos(6t) + e^{3t}\sin(4t)$\\
 \item  $\ds f(t) = t^2 \sin(3t)$\\
 \een
%\item  \ben [label=(\alph*)]
% \item  $\ds F(s) = \frac{2}{(s+2)^2 + 4} + \frac{2}{(s-3)^3}, \hspace{0.25cm}s>3$\\
% \item  $\ds G(s) = \frac{6}{(s+9)^4}, \hspace{0.25cm}s>-9$\\
% \item  $\ds H(s) = \frac{1}{s-2} - \frac{6}{s^4} + \frac{2}{s^3} - \frac{5}{s^2 + 25},\hspace{0.25cm}s>2$\\
% \item  $\ds Q(s) = \frac{8(s^2-36)}{(s^2 + 36)^2} + \frac{4}{(s-3)^2 + 16}, \hspace{0.25cm}s>3$\\
% \item  $\ds F(s) = \frac{18(s^2 - 3)}{(s^2+9)^3},\hspace{0.25cm}s>0$\\
% \een
\end{DEPP}
%%%%%%%%%%%%%%%%%%%%%%%%%%%%%%%%%%%%%%%%%%%%%%%%%%%%%
%%%%%%%%%%%%%%%%%%%%%%%%%%%%%%%%%%%%%%%%%%%%%%%%%%%%%
\showto{answers}{
\newpage
\noindent  Answers:
\begin{DEPP}
\item    $\ds H(s) = \frac{7}{s^2},$ $s>0$\\

\item    $\ds W(s) = \frac{1}{s-k},$  $s > k$ \\

\item   $\ds Y(s) = \frac{2}{s^2} + \frac{1}{s-5},$  $s>5$\\

\item  $\ds D(s) = \frac{7}{s} - \frac{7e^{-3s}}{s},$  $s>0$\\

\item  \ben [label=(\alph*)]
 \item  $\ds F(s) = \frac{2}{(s+2)^2 + 4} + \frac{2}{(s-3)^3}, \hspace{0.25cm}s>3$\\
 \item  $\ds G(s) = \frac{6}{(s+9)^4}, \hspace{0.25cm}s>-9$\\
 \item  $\ds H(s) = \frac{1}{s-2} - \frac{6}{s^4} + \frac{2}{s^3} - \frac{5}{s^2 + 25},\hspace{0.25cm}s>2$\\
 \item  $\ds Q(s) = \frac{8(s^2-36)}{(s^2 + 36)^2} + \frac{4}{(s-3)^2 + 16}, \hspace{0.25cm}s>3$\\
 \item  $\ds F(s) = \frac{18(s^2 - 3)}{(s^2+9)^3},\hspace{0.25cm}s>0$\\
 \een
\end{DEPP}
}


%%%%%%%%%%%%%%%%%%%%%%%%%%%%%%%%%%%%%%%%%%%%%%%%%%%%%
\newpage

%\mbox{}\newpage  %If needed

\showto{answers}{\mbox{}\hfill ANSWERS\\}
\noindent\rule{7in}{0.4pt}\\

{\sc
\noindent Differential Equations       \showto{student}{\hfill Name:  \line(1,0){200}}\\
DEs Performance Quiz  \DEPPSetNum\   \showto{student}{\hfill HR:  \line(1,0){200}}\\
Due by 1600 \PQDueDate \showto{student}{\hfill Circle your section: \hspace{0.73in}\mbox{}\\
              \mbox{} \hfill MWF1000
               \hspace{0.5cm}MWF1300 \hspace{0.5cm}TR0925
               \hspace{0.5cm}TR1050}\\
}

\noindent\rule{7in}{0.4pt}
\vspace{0.5cm}

\noindent  For this Performance Quiz, you are permitted to use your own notes/work and your
textbook, but you are NOT AUTHORIZED to seek help from anyone
other than MAJ Libertini or MAJ Bliss.\\

%%%%%%%%%%%%%%%%%%%%%%%%%%%%%%%%%%%%%%%%%%%%%%%%%%%%%
\begin{PQ}
\item   {\bf Use the definition of Laplace Transform} to find the Laplace Transform of

  \[y(t) = \left\{ \begin{array}{ll} 0, & t<4\\
        6, & t \ge 4\\
       \end{array} \right.\]    %=6U(t-4)

  Include appropriate restrictions on $s$ to guarantee the convergence of the integral.\\

  {\it Note:  Credit will only be assigned for work that demonstrates how the definition was
  used to find the LT.}\\
\vfill
\begin{center}{\it (There's more on the back!)}
\end{center}
\newpage
\item  Use the table and properties provided to determine the Laplace transform of each function.
Include resulting restrictions on $s.$
 \ben[label=(\alph*)]
 \item  $\ds g(t) = 17t^3 - 9te^{-4t}$\\ \vfill
 \item  $\ds h(t) = t^2 \cos(8t)$\\ \vfill
 \een
\end{PQ}
%%%%%%%%%%%%%%%%%%%%%%%%%%%%%%%%%%%%%%%%%%%%%%%%%%%%%
%%%%%%%%%%%%%%%%%%%%%%%%%%%%%%%%%%%%%%%%%%%%%%%%%%%%%
\showto{answers}{
\newpage
\noindent  Answers:
\begin{PQ}
\item  $\ds Y(s) = \frac{6e^{-4s}}{s},$  $s>0$

\item  \ben[label=(\alph*)]
 \item  $\ds G(s) = \frac{102}{s^4} - \frac{9}{(s+4)^2},$ $s>0$\\

 \item  $\ds H(s) = -16 \cdot \frac{64-3s^2}{(s^2+64)^3}$
  %H(s) = (-1)^2 *d^2/ds^2[8/(s^2 + 64)]
  %  = d/ds[( (s^2+64)(0) - 8(2s)) / (s^2+64)^2]
  %  = -16*d/ds[s / (s^2+64)^2]
  %  = -16 * [ ( (s^2+64)^2 - s*2(s^2+64)(2s) ) / (s^2+64)^4]
  %  = -16 * [ ( (s^2 + 64) - 4s^2 ) / (s^2+64)^3]
  %  = -16 * [ ( -3s^2 + 64) / (s^2+64)^3]
 \een
\end{PQ}
}
%%%%%%%%%%%%%%%%%%%%%%%%%%%%%%%%%%%%%%%%%%%%%%%%%%%%%
%%%%%%%%%%%%%%%%%%%%%%%%%%%%%%%%%%%%%%%%%%%%%%%%%%%%%

\fi

%\ifnum\conditionmacro=3 
%	\def\CurrentAudience{solutions}						%EDIT (4) HERE AS NEEDED
%	\pagestyle{empty}

\showto{answers}{\hfill ANSWERS\\}
\noindent\rule{7in}{0.4pt}
\begin{center}
{\sc
\noindent Differential Equations\\
 DEs Practice Problems Set \DEPPSetNum\        %CHANGE NUMBER BELOW********
}
\end{center}
\noindent\rule{7in}{0.4pt}

\noindent Submit solutions for extra credit by 1600 \ECDueDate\\
\noindent Corresponding Performance Quiz due by 1600 \PQDueDate\\

\vspace{0.5cm}

%intro to LTs / basic LTs.
%%%%%%%%%%%%%%%%%%%%%%%%%%%%%%%%%%%%%%%%%%%%%%%%%%%%%
\subsection*{The Definition of Laplace Transform}

\noindent  As a reminder, the Laplace Transform of function $f(t)$ is a function $F(s)$ as follows:
 \[ F(s) = \L\{ f(t) \} = \int_0^{\infty}e^{-st}f(t) dt, \]
 provided this integral is finite.\\

\begin{DEPP}
\item    Use the definition of Laplace Transform (above) to find the Laplace Transform of $h(t) = 7t.$
 Include appropriate restrictions on $s$ to guarantee the convergence of the integral.\\

\item    Use the definition of Laplace Transform (above) to find the Laplace Transform of $w(t) = e^{kt},$
 where $k$ is some constant.  Include appropriate restrictions on $s$ to guarantee the
 convergence of the integral.\\

\item   Use the definition of Laplace Transform (above) to find the Laplace Transform of $y(t) = 7t + e^{5t}.$
  Include appropriate restrictions on $s$ to guarantee the convergence of the integral.\\

\item  Use the definition of Laplace Transform (above) to find the Laplace Transform of  the function
 \[d(t) = \left\{ \begin{array}{ll} 0, & t<0\\
       7, & 0 \le t < 3\\
       0, & t \ge 3\\
       \end{array} \right.\]  %= 7U(t) - 7U(t-3)
 Include appropriate restrictions on $s$ to guarantee the convergence of the integral.\\

 {\it Note:  This is a really important skill!  If we can take the Laplace Transform of
 piecewise-defined functions like this one, then we can solve differential equations that
 we couldn't solve before!)}\\
%\item    $\ds H(s) = \frac{7}{s^2},$ $s>0$\\
%
%\item    $\ds W(s) = \frac{1}{s-k},$  $s > k$ \\
%
%\item   $\ds Y(s) = \frac{2}{s^2} + \frac{1}{s-5},$  $s>5$

%\item  $\ds D(s) = \frac{7}{s} - \frac{7e^{-3s}}{s}$
\end{DEPP}
\newpage
\subsection*{Using Table to Find LTs}
\begin{DEPP}[resume]
\item  Use the table and properties provided to determine the Laplace transform of each function.
Include resulting restrictions on $s.$
 \ben[label=(\alph*)]
 \item  $\ds f(t) = e^{-2t}\sin(2t) + t^2 e^{3t}$\\
 \item  $\ds g(t) = t^3 e^{-9t}$\\
 \item  $\ds h(t) = e^{2t} - t^3 + t^2 - \sin (5t)$\\
 \item  $\ds q(t) = 8t\cos(6t) + e^{3t}\sin(4t)$\\
 \item  $\ds f(t) = t^2 \sin(3t)$\\
 \een
%\item  \ben [label=(\alph*)]
% \item  $\ds F(s) = \frac{2}{(s+2)^2 + 4} + \frac{2}{(s-3)^3}, \hspace{0.25cm}s>3$\\
% \item  $\ds G(s) = \frac{6}{(s+9)^4}, \hspace{0.25cm}s>-9$\\
% \item  $\ds H(s) = \frac{1}{s-2} - \frac{6}{s^4} + \frac{2}{s^3} - \frac{5}{s^2 + 25},\hspace{0.25cm}s>2$\\
% \item  $\ds Q(s) = \frac{8(s^2-36)}{(s^2 + 36)^2} + \frac{4}{(s-3)^2 + 16}, \hspace{0.25cm}s>3$\\
% \item  $\ds F(s) = \frac{18(s^2 - 3)}{(s^2+9)^3},\hspace{0.25cm}s>0$\\
% \een
\end{DEPP}
%%%%%%%%%%%%%%%%%%%%%%%%%%%%%%%%%%%%%%%%%%%%%%%%%%%%%
%%%%%%%%%%%%%%%%%%%%%%%%%%%%%%%%%%%%%%%%%%%%%%%%%%%%%
\showto{answers}{
\newpage
\noindent  Answers:
\begin{DEPP}
\item    $\ds H(s) = \frac{7}{s^2},$ $s>0$\\

\item    $\ds W(s) = \frac{1}{s-k},$  $s > k$ \\

\item   $\ds Y(s) = \frac{2}{s^2} + \frac{1}{s-5},$  $s>5$\\

\item  $\ds D(s) = \frac{7}{s} - \frac{7e^{-3s}}{s},$  $s>0$\\

\item  \ben [label=(\alph*)]
 \item  $\ds F(s) = \frac{2}{(s+2)^2 + 4} + \frac{2}{(s-3)^3}, \hspace{0.25cm}s>3$\\
 \item  $\ds G(s) = \frac{6}{(s+9)^4}, \hspace{0.25cm}s>-9$\\
 \item  $\ds H(s) = \frac{1}{s-2} - \frac{6}{s^4} + \frac{2}{s^3} - \frac{5}{s^2 + 25},\hspace{0.25cm}s>2$\\
 \item  $\ds Q(s) = \frac{8(s^2-36)}{(s^2 + 36)^2} + \frac{4}{(s-3)^2 + 16}, \hspace{0.25cm}s>3$\\
 \item  $\ds F(s) = \frac{18(s^2 - 3)}{(s^2+9)^3},\hspace{0.25cm}s>0$\\
 \een
\end{DEPP}
}


%%%%%%%%%%%%%%%%%%%%%%%%%%%%%%%%%%%%%%%%%%%%%%%%%%%%%
\newpage

%\mbox{}\newpage  %If needed

\showto{answers}{\mbox{}\hfill ANSWERS\\}
\noindent\rule{7in}{0.4pt}\\

{\sc
\noindent Differential Equations       \showto{student}{\hfill Name:  \line(1,0){200}}\\
DEs Performance Quiz  \DEPPSetNum\   \showto{student}{\hfill HR:  \line(1,0){200}}\\
Due by 1600 \PQDueDate \showto{student}{\hfill Circle your section: \hspace{0.73in}\mbox{}\\
              \mbox{} \hfill MWF1000
               \hspace{0.5cm}MWF1300 \hspace{0.5cm}TR0925
               \hspace{0.5cm}TR1050}\\
}

\noindent\rule{7in}{0.4pt}
\vspace{0.5cm}

\noindent  For this Performance Quiz, you are permitted to use your own notes/work and your
textbook, but you are NOT AUTHORIZED to seek help from anyone
other than MAJ Libertini or MAJ Bliss.\\

%%%%%%%%%%%%%%%%%%%%%%%%%%%%%%%%%%%%%%%%%%%%%%%%%%%%%
\begin{PQ}
\item   {\bf Use the definition of Laplace Transform} to find the Laplace Transform of

  \[y(t) = \left\{ \begin{array}{ll} 0, & t<4\\
        6, & t \ge 4\\
       \end{array} \right.\]    %=6U(t-4)

  Include appropriate restrictions on $s$ to guarantee the convergence of the integral.\\

  {\it Note:  Credit will only be assigned for work that demonstrates how the definition was
  used to find the LT.}\\
\vfill
\begin{center}{\it (There's more on the back!)}
\end{center}
\newpage
\item  Use the table and properties provided to determine the Laplace transform of each function.
Include resulting restrictions on $s.$
 \ben[label=(\alph*)]
 \item  $\ds g(t) = 17t^3 - 9te^{-4t}$\\ \vfill
 \item  $\ds h(t) = t^2 \cos(8t)$\\ \vfill
 \een
\end{PQ}
%%%%%%%%%%%%%%%%%%%%%%%%%%%%%%%%%%%%%%%%%%%%%%%%%%%%%
%%%%%%%%%%%%%%%%%%%%%%%%%%%%%%%%%%%%%%%%%%%%%%%%%%%%%
\showto{answers}{
\newpage
\noindent  Answers:
\begin{PQ}
\item  $\ds Y(s) = \frac{6e^{-4s}}{s},$  $s>0$

\item  \ben[label=(\alph*)]
 \item  $\ds G(s) = \frac{102}{s^4} - \frac{9}{(s+4)^2},$ $s>0$\\

 \item  $\ds H(s) = -16 \cdot \frac{64-3s^2}{(s^2+64)^3}$
  %H(s) = (-1)^2 *d^2/ds^2[8/(s^2 + 64)]
  %  = d/ds[( (s^2+64)(0) - 8(2s)) / (s^2+64)^2]
  %  = -16*d/ds[s / (s^2+64)^2]
  %  = -16 * [ ( (s^2+64)^2 - s*2(s^2+64)(2s) ) / (s^2+64)^4]
  %  = -16 * [ ( (s^2 + 64) - 4s^2 ) / (s^2+64)^3]
  %  = -16 * [ ( -3s^2 + 64) / (s^2+64)^3]
 \een
\end{PQ}
}
%%%%%%%%%%%%%%%%%%%%%%%%%%%%%%%%%%%%%%%%%%%%%%%%%%%%%
%%%%%%%%%%%%%%%%%%%%%%%%%%%%%%%%%%%%%%%%%%%%%%%%%%%%%

%\fi

\end{document}
